%------------------------------------------------
%  Project: CRLM_dormancy
%  Report: GSE114012 dormancy analysis
%  Theme: SimplePlus Beamer
%------------------------------------------------
\makeatletter
\def\input@path{{/Users/solomon/Projects/CRLM_dormancy/data/templates/beamer/SimplePlus-BeamerTheme/}}
\makeatother

\documentclass[11pt,aspectratio=169,xcolor=dvipsnames]{beamer}
\usetheme{SimplePlus}
\usepackage[UTF8]{ctex}
\usepackage{graphicx}
\usepackage{booktabs}
\usepackage{array}
\usepackage{adjustbox}
\usepackage{fancyvrb}
\usepackage{hyperref}
\usepackage{bookmark}
\usepackage{url}
\usepackage{xcolor}
\setbeamertemplate{navigation symbols}{}
\setbeamersize{text margin left=1.5mm,text margin right=1.5mm}
\setbeamerfont{normal text}{size=\scriptsize}
\setbeamerfont{itemize/enumerate body}{size=\scriptsize}
\setbeamerfont{itemize/enumerate subbody}{size=\tiny}
\AtBeginEnvironment{frame}{\fontsize{7.2pt}{8.4pt}\selectfont}
\hfuzz=5pt
\vfuzz=10pt
\hbadness=10000
\vbadness=10000


\newcommand{\CRLMROOT}{/Users/solomon/Projects/CRLM_dormancy}
\newcommand{\CRLMPath}[1]{\begingroup\tiny\nolinkurl{#1}\endgroup}
\newsavebox{\CRLMFigureBox}
\newlength{\CRLMFigureTextWidth}

\newcommand{\CRLMMissingFigure}[1]{%
  \fbox{%
    \begin{minipage}[c][0.82\textheight][c]{0.96\linewidth}
      \centering
      \ttfamily\scriptsize
      \textbf{Missing figure}\\
      \texttt{#1}\\
      \texttt{please check path and output}
    \end{minipage}%
  }%
}
\newcommand{\FrameFig}[4]{%
  \begin{frame}[plain]
    \vspace{0.15em}
    {\scriptsize\bfseries #1}
    \par\vspace{0.18em}
    \IfFileExists{\CRLMROOT/#2}{%
      \sbox{\CRLMFigureBox}{%
        \includegraphics[
          width=0.72\textwidth,
          height=0.87\paperheight,
          keepaspectratio
        ]{\CRLMROOT/#2}%
      }%
      \setlength{\CRLMFigureTextWidth}{\dimexpr\textwidth-\wd\CRLMFigureBox\relax}%
      \noindent
      \begin{adjustbox}{valign=t}%
        \usebox{\CRLMFigureBox}%
      \end{adjustbox}%
      \begin{adjustbox}{valign=t}%
        \begin{minipage}{\CRLMFigureTextWidth}%
          \fontsize{7.2pt}{8.6pt}\selectfont
          \textbf{说明：}
          \begin{itemize}
            \setlength{\itemsep}{0.25em}
            \setlength{\leftmargini}{1.1em}
            #4
          \end{itemize}
          \vspace{0.45em}
          \textbf{相对路径：}\par
          \CRLMPath{#2}
        \end{minipage}%
      \end{adjustbox}%
    }{%
      \CRLMMissingFigure{#2}%
    }%
  \end{frame}
}
\newcommand{\FrameTable}[4]{%
  \begin{frame}{#1}
    \vspace{-0.35em}
    \scriptsize
    \textbf{说明：} #3
    \par\vspace{0.18em}
    \textbf{相对路径：} \CRLMPath{#2}
    \par\vspace{0.18em}
    \textbf{表格预览：} 展示全部列与前21行。
    \par\vspace{0.22em}
    \begin{adjustbox}{max width=\textwidth,max totalheight=0.68\textheight,keepaspectratio}
      \input{\CRLMROOT/#4}
    \end{adjustbox}
  \end{frame}
}

\title[GSE114012 休眠细胞复苏]{结直肠癌休眠细胞复苏的转录组线索}
\subtitle{ }
\author{王瀚 郭鹏程}
\institute{ }
\date{\today}

\begin{document}

\begin{frame}
  \titlepage
\end{frame}

\begin{frame}{研究背景与核心假设}
  \begin{itemize}
    \item 结直肠癌远处复发可能并非单纯由晚期转移扩散决定，部分播散或残留癌细胞可能在早期已进入远处生态位。
    \item 本课题关注休眠癌细胞从低增殖状态重新苏醒的分子触发因素，尤其是转录调控、通路活化与微环境信号。
    \item 核心假说：若能识别并阻断促使休眠细胞复苏的关键轴，就可能使残留癌细胞长期维持休眠，从而降低复发与转移风险。
    \item 本页汇报先以 GSE114012 的 LRC/BULK RNA-seq 模型建立可复现的候选基因证据链，为后续 TF、通路和药物重定位分析提供入口。
  \end{itemize}
\end{frame}

\begin{frame}{GSE114012 数据模型与本次分析任务}
  \begin{itemize}
    \item GEO 数据集 \texttt{GSE114012} 的主题是识别结直肠癌球体培养中的休眠细胞；官方设计使用 CFSE 标记并分选保留染料的 LRC 与增殖/循环状态的 BULK 细胞。
    \item 数据包含 DLD1、HCT15、HT55、SW948、RKO、SW48 六种结直肠癌细胞系，每种细胞系包含 LRC 与 BULK 群体的重复样本。
    \item 本次分析将 LRC 视作休眠样细胞状态、BULK 视作相对循环/增殖状态，重点比较 LRC vs BULK 的表达差异与跨细胞系稳定交集。
    \item 所有差异分析、交集和火山图共享同一显著性阈值，以保证后续候选“休眠维持/复苏相关基因”的可比性。
  \end{itemize}
\end{frame}

\begin{frame}{脚本流程与结果约束}
  \begin{enumerate}
    \item \textbf{00}：使用 TPM 做 LRC 与 BULK 样本层级聚类，先检查休眠样/增殖样样本结构与离群情况。
    \item \textbf{01}：基于 limma-voom 进行 LRC vs BULK 差异分析，得到每个细胞系及组合设计的显著基因。
    \item \textbf{02}：按可配置细胞系组合提取显著基因交集，筛选跨模型重复出现的候选休眠相关信号。
    \item \textbf{03}：用传统火山图观察单个比较中的表达漂移、显著性强度和候选基因位置。
    \item \textbf{04}：用多组火山图并列比较多套细胞系组合，突出方向一致的上调/下调信号。
    \item \textbf{05}：提取每套差异分析设计的 Top 50 显著差异基因，并用 TPM 行 Z-Score 热图展示 LRC/BULK 表达模式。
    \item \textbf{06}：基于完整差异分析排序结果批量运行 MSigDB GSEA，并用 dotplot 与单通路 GSEA+热图解释休眠样状态相关通路。
  \end{enumerate}
\end{frame}

\FrameFig
  {00. LRC 样本热图（TPM）}
  {results/ngs/GSE114012/plots/sample_clustering_heatmap/coding/LRC/heatmap.pdf}
  {LRC 全样本聚类展示，保留样本标题（Title）用于标签。}
  {
    \item LRC 是 CFSE 保留的休眠样群体，聚类用于观察不同细胞系休眠样样本是否呈现相近表达结构；
    \item 若 LRC 样本之间存在稳定聚类，说明休眠样状态可能具有跨模型共享的转录轮廓；
    \item 该图为后续差异分析提供样本质量与分组结构依据。
  }

\FrameFig
  {00. BULK 样本热图（TPM）}
  {results/ngs/GSE114012/plots/sample_clustering_heatmap/coding/BULK/heatmap.pdf}
  {BULK 全样本聚类展示。}
  {
    \item BULK 代表相对循环/增殖状态，用于与 LRC 休眠样群体建立背景对照；
    \item 聚类结构可帮助判断 BULK 样本是否主要受细胞系来源或增殖状态驱动；
    \item 与 LRC 热图对照可辅助解释休眠样表达差异是否具有模型一致性。
  }

\begin{frame}{01. 差异分析目录与统一输出结构}
  \small
  每个分组分析都输出如下固定三张表，放在统一目录：
  \vspace{0.22em}
  \texttt{results/ngs/GSE114012/tables/<Analysis\_Name>/DEG/}
  \vspace{0.22em}
  \begin{itemize}
    \item \texttt{all\_genes.csv}
    \item \texttt{significant\_genes.csv}
    \item \texttt{summary.csv}
  \end{itemize}
  \vspace{0.22em}
  下方所有“单张表页”按同一风格逐一展示，默认给出文件前若干行以便快速核对字段与结果质量。
\end{frame}

\FrameTable
  {01. 差异分析 Summary：ALL}
  {results/ngs/GSE114012/tables/ALL/DEG/summary.csv}
  {该表概括当前 LRC vs BULK 比较的样本数、显著上调/下调基因数和阈值设置，用于判断每个细胞系或组合模型中休眠样 LRC 相对 BULK 的总体转录变化强度。}
  {results/ngs/GSE114012/tables/ALL/DEG/summary.tex}


\FrameTable
  {01. ALL：all\_genes}
  {results/ngs/GSE114012/tables/ALL/DEG/all_genes.csv}
  {该表保留完整差异分析排序结果，可用于追溯每个基因在 LRC/BULK 模型中的 logFC、显著性和排序位置，是后续候选 TF、通路和药物重定位输入的基础。}
  {results/ngs/GSE114012/tables/ALL/DEG/all_genes.tex}


\FrameTable
  {01. ALL：significant\_genes}
  {results/ngs/GSE114012/tables/ALL/DEG/significant_genes.csv}
  {该表按统一阈值筛选显著差异基因，代表当前模型中最直接的休眠样 LRC 转录特征，并作为交集分析和火山图展示的核心输入。}
  {results/ngs/GSE114012/tables/ALL/DEG/significant_genes.tex}


\FrameTable
  {01. 差异分析 Summary：DLD1}
  {results/ngs/GSE114012/tables/DLD1/DEG/summary.csv}
  {该表概括当前 LRC vs BULK 比较的样本数、显著上调/下调基因数和阈值设置，用于判断每个细胞系或组合模型中休眠样 LRC 相对 BULK 的总体转录变化强度。}
  {results/ngs/GSE114012/tables/DLD1/DEG/summary.tex}


\FrameTable
  {01. DLD1：all\_genes}
  {results/ngs/GSE114012/tables/DLD1/DEG/all_genes.csv}
  {该表保留完整差异分析排序结果，可用于追溯每个基因在 LRC/BULK 模型中的 logFC、显著性和排序位置，是后续候选 TF、通路和药物重定位输入的基础。}
  {results/ngs/GSE114012/tables/DLD1/DEG/all_genes.tex}


\FrameTable
  {01. DLD1：significant\_genes}
  {results/ngs/GSE114012/tables/DLD1/DEG/significant_genes.csv}
  {该表按统一阈值筛选显著差异基因，代表当前模型中最直接的休眠样 LRC 转录特征，并作为交集分析和火山图展示的核心输入。}
  {results/ngs/GSE114012/tables/DLD1/DEG/significant_genes.tex}


\FrameTable
  {01. 差异分析 Summary：DLD1/HCT15}
  {results/ngs/GSE114012/tables/DLD1_HCT15/DEG/summary.csv}
  {该表概括当前 LRC vs BULK 比较的样本数、显著上调/下调基因数和阈值设置，用于判断每个细胞系或组合模型中休眠样 LRC 相对 BULK 的总体转录变化强度。}
  {results/ngs/GSE114012/tables/DLD1_HCT15/DEG/summary.tex}


\FrameTable
  {01. DLD1/HCT15：all\_genes}
  {results/ngs/GSE114012/tables/DLD1_HCT15/DEG/all_genes.csv}
  {该表保留完整差异分析排序结果，可用于追溯每个基因在 LRC/BULK 模型中的 logFC、显著性和排序位置，是后续候选 TF、通路和药物重定位输入的基础。}
  {results/ngs/GSE114012/tables/DLD1_HCT15/DEG/all_genes.tex}


\FrameTable
  {01. DLD1/HCT15：significant\_genes}
  {results/ngs/GSE114012/tables/DLD1_HCT15/DEG/significant_genes.csv}
  {该表按统一阈值筛选显著差异基因，代表当前模型中最直接的休眠样 LRC 转录特征，并作为交集分析和火山图展示的核心输入。}
  {results/ngs/GSE114012/tables/DLD1_HCT15/DEG/significant_genes.tex}


\FrameTable
  {01. 差异分析 Summary：DLD1/HCT15/SW48}
  {results/ngs/GSE114012/tables/DLD1_HCT15_SW48/DEG/summary.csv}
  {该表概括当前 LRC vs BULK 比较的样本数、显著上调/下调基因数和阈值设置，用于判断每个细胞系或组合模型中休眠样 LRC 相对 BULK 的总体转录变化强度。}
  {results/ngs/GSE114012/tables/DLD1_HCT15_SW48/DEG/summary.tex}


\FrameTable
  {01. DLD1/HCT15/SW48：all\_genes}
  {results/ngs/GSE114012/tables/DLD1_HCT15_SW48/DEG/all_genes.csv}
  {该表保留完整差异分析排序结果，可用于追溯每个基因在 LRC/BULK 模型中的 logFC、显著性和排序位置，是后续候选 TF、通路和药物重定位输入的基础。}
  {results/ngs/GSE114012/tables/DLD1_HCT15_SW48/DEG/all_genes.tex}


\FrameTable
  {01. DLD1/HCT15/SW48：significant\_genes}
  {results/ngs/GSE114012/tables/DLD1_HCT15_SW48/DEG/significant_genes.csv}
  {该表按统一阈值筛选显著差异基因，代表当前模型中最直接的休眠样 LRC 转录特征，并作为交集分析和火山图展示的核心输入。}
  {results/ngs/GSE114012/tables/DLD1_HCT15_SW48/DEG/significant_genes.tex}


\FrameTable
  {01. 差异分析 Summary：HCT15}
  {results/ngs/GSE114012/tables/HCT15/DEG/summary.csv}
  {该表概括当前 LRC vs BULK 比较的样本数、显著上调/下调基因数和阈值设置，用于判断每个细胞系或组合模型中休眠样 LRC 相对 BULK 的总体转录变化强度。}
  {results/ngs/GSE114012/tables/HCT15/DEG/summary.tex}


\FrameTable
  {01. HCT15：all\_genes}
  {results/ngs/GSE114012/tables/HCT15/DEG/all_genes.csv}
  {该表保留完整差异分析排序结果，可用于追溯每个基因在 LRC/BULK 模型中的 logFC、显著性和排序位置，是后续候选 TF、通路和药物重定位输入的基础。}
  {results/ngs/GSE114012/tables/HCT15/DEG/all_genes.tex}


\FrameTable
  {01. HCT15：significant\_genes}
  {results/ngs/GSE114012/tables/HCT15/DEG/significant_genes.csv}
  {该表按统一阈值筛选显著差异基因，代表当前模型中最直接的休眠样 LRC 转录特征，并作为交集分析和火山图展示的核心输入。}
  {results/ngs/GSE114012/tables/HCT15/DEG/significant_genes.tex}


\FrameTable
  {01. 差异分析 Summary：HT55}
  {results/ngs/GSE114012/tables/HT55/DEG/summary.csv}
  {该表概括当前 LRC vs BULK 比较的样本数、显著上调/下调基因数和阈值设置，用于判断每个细胞系或组合模型中休眠样 LRC 相对 BULK 的总体转录变化强度。}
  {results/ngs/GSE114012/tables/HT55/DEG/summary.tex}


\FrameTable
  {01. HT55：all\_genes}
  {results/ngs/GSE114012/tables/HT55/DEG/all_genes.csv}
  {该表保留完整差异分析排序结果，可用于追溯每个基因在 LRC/BULK 模型中的 logFC、显著性和排序位置，是后续候选 TF、通路和药物重定位输入的基础。}
  {results/ngs/GSE114012/tables/HT55/DEG/all_genes.tex}


\FrameTable
  {01. HT55：significant\_genes}
  {results/ngs/GSE114012/tables/HT55/DEG/significant_genes.csv}
  {该表按统一阈值筛选显著差异基因，代表当前模型中最直接的休眠样 LRC 转录特征，并作为交集分析和火山图展示的核心输入。}
  {results/ngs/GSE114012/tables/HT55/DEG/significant_genes.tex}


\FrameTable
  {01. 差异分析 Summary：RKO}
  {results/ngs/GSE114012/tables/RKO/DEG/summary.csv}
  {该表概括当前 LRC vs BULK 比较的样本数、显著上调/下调基因数和阈值设置，用于判断每个细胞系或组合模型中休眠样 LRC 相对 BULK 的总体转录变化强度。}
  {results/ngs/GSE114012/tables/RKO/DEG/summary.tex}


\FrameTable
  {01. RKO：all\_genes}
  {results/ngs/GSE114012/tables/RKO/DEG/all_genes.csv}
  {该表保留完整差异分析排序结果，可用于追溯每个基因在 LRC/BULK 模型中的 logFC、显著性和排序位置，是后续候选 TF、通路和药物重定位输入的基础。}
  {results/ngs/GSE114012/tables/RKO/DEG/all_genes.tex}


\FrameTable
  {01. RKO：significant\_genes}
  {results/ngs/GSE114012/tables/RKO/DEG/significant_genes.csv}
  {该表按统一阈值筛选显著差异基因，代表当前模型中最直接的休眠样 LRC 转录特征，并作为交集分析和火山图展示的核心输入。}
  {results/ngs/GSE114012/tables/RKO/DEG/significant_genes.tex}


\FrameTable
  {01. 差异分析 Summary：SW48}
  {results/ngs/GSE114012/tables/SW48/DEG/summary.csv}
  {该表概括当前 LRC vs BULK 比较的样本数、显著上调/下调基因数和阈值设置，用于判断每个细胞系或组合模型中休眠样 LRC 相对 BULK 的总体转录变化强度。}
  {results/ngs/GSE114012/tables/SW48/DEG/summary.tex}


\FrameTable
  {01. SW48：all\_genes}
  {results/ngs/GSE114012/tables/SW48/DEG/all_genes.csv}
  {该表保留完整差异分析排序结果，可用于追溯每个基因在 LRC/BULK 模型中的 logFC、显著性和排序位置，是后续候选 TF、通路和药物重定位输入的基础。}
  {results/ngs/GSE114012/tables/SW48/DEG/all_genes.tex}


\FrameTable
  {01. SW48：significant\_genes}
  {results/ngs/GSE114012/tables/SW48/DEG/significant_genes.csv}
  {该表按统一阈值筛选显著差异基因，代表当前模型中最直接的休眠样 LRC 转录特征，并作为交集分析和火山图展示的核心输入。}
  {results/ngs/GSE114012/tables/SW48/DEG/significant_genes.tex}


\FrameTable
  {01. 差异分析 Summary：SW948}
  {results/ngs/GSE114012/tables/SW948/DEG/summary.csv}
  {该表概括当前 LRC vs BULK 比较的样本数、显著上调/下调基因数和阈值设置，用于判断每个细胞系或组合模型中休眠样 LRC 相对 BULK 的总体转录变化强度。}
  {results/ngs/GSE114012/tables/SW948/DEG/summary.tex}


\FrameTable
  {01. SW948：all\_genes}
  {results/ngs/GSE114012/tables/SW948/DEG/all_genes.csv}
  {该表保留完整差异分析排序结果，可用于追溯每个基因在 LRC/BULK 模型中的 logFC、显著性和排序位置，是后续候选 TF、通路和药物重定位输入的基础。}
  {results/ngs/GSE114012/tables/SW948/DEG/all_genes.tex}


\FrameTable
  {01. SW948：significant\_genes}
  {results/ngs/GSE114012/tables/SW948/DEG/significant_genes.csv}
  {该表按统一阈值筛选显著差异基因，代表当前模型中最直接的休眠样 LRC 转录特征，并作为交集分析和火山图展示的核心输入。}
  {results/ngs/GSE114012/tables/SW948/DEG/significant_genes.tex}


\begin{frame}{02. 交集分析说明}
  \begin{itemize}
    \item 输出目录：\texttt{results/ngs/GSE114012/intersect/<scheme>/}
  \end{itemize}
  \vspace{0.22em}
  \begin{itemize}
    \item \texttt{summary.csv}：每个交集方案的统计汇总（交集基因数、共同上调、共同下调等）。
    \item \texttt{gene\_list.csv}：交集基因及其注释（不含模型统计量列）。
    \item \texttt{deg\_results.csv}：交集基因在方案内每个参与分析中的完整 DEG 行信息。
  \end{itemize}
\end{frame}

\FrameTable
  {02. DLD1/HCT15 交集统计}
  {results/ngs/GSE114012/intersect/DLD1_HCT15/summary.csv}
  {该表汇总当前交集方案中跨细胞系重复出现的显著基因数量及共同上调/下调方向，用于优先筛选稳定的休眠维持或复苏相关候选信号。}
  {results/ngs/GSE114012/intersect/DLD1_HCT15/summary.tex}


\FrameTable
  {02. DLD1/HCT15 交集基因列表}
  {results/ngs/GSE114012/intersect/DLD1_HCT15/gene_list.csv}
  {该表列出交集后的基因清单和基础注释，去除了模型统计量，便于将稳定候选基因继续输入 TF 富集、通路富集和药物重定位分析。}
  {results/ngs/GSE114012/intersect/DLD1_HCT15/gene_list.tex}


\FrameTable
  {02. DLD1/HCT15 方案中的 DLD1 DEG 子集}
  {results/ngs/GSE114012/intersect/DLD1_HCT15/DLD1/deg_results.csv}
  {该表将交集基因回填到对应差异分析结果中，展示同一候选基因在单个细胞系 LRC vs BULK 比较中的方向、效应量和显著性。}
  {results/ngs/GSE114012/intersect/DLD1_HCT15/DLD1/deg_results.tex}


\FrameTable
  {02. DLD1/HCT15 方案中的 HCT15 DEG 子集}
  {results/ngs/GSE114012/intersect/DLD1_HCT15/HCT15/deg_results.csv}
  {该表将交集基因回填到对应差异分析结果中，展示同一候选基因在单个细胞系 LRC vs BULK 比较中的方向、效应量和显著性。}
  {results/ngs/GSE114012/intersect/DLD1_HCT15/HCT15/deg_results.tex}


\FrameTable
  {02. DLD1/HCT15/SW48 交集统计}
  {results/ngs/GSE114012/intersect/DLD1_HCT15_SW48/summary.csv}
  {该表汇总当前交集方案中跨细胞系重复出现的显著基因数量及共同上调/下调方向，用于优先筛选稳定的休眠维持或复苏相关候选信号。}
  {results/ngs/GSE114012/intersect/DLD1_HCT15_SW48/summary.tex}


\FrameTable
  {02. DLD1/HCT15/SW48 交集基因列表}
  {results/ngs/GSE114012/intersect/DLD1_HCT15_SW48/gene_list.csv}
  {该表列出交集后的基因清单和基础注释，去除了模型统计量，便于将稳定候选基因继续输入 TF 富集、通路富集和药物重定位分析。}
  {results/ngs/GSE114012/intersect/DLD1_HCT15_SW48/gene_list.tex}


\FrameTable
  {02. DLD1/HCT15/SW48 方案中的 DLD1 DEG 子集}
  {results/ngs/GSE114012/intersect/DLD1_HCT15_SW48/DLD1/deg_results.csv}
  {该表将交集基因回填到对应差异分析结果中，展示同一候选基因在单个细胞系 LRC vs BULK 比较中的方向、效应量和显著性。}
  {results/ngs/GSE114012/intersect/DLD1_HCT15_SW48/DLD1/deg_results.tex}


\FrameTable
  {02. DLD1/HCT15/SW48 方案中的 HCT15 DEG 子集}
  {results/ngs/GSE114012/intersect/DLD1_HCT15_SW48/HCT15/deg_results.csv}
  {该表将交集基因回填到对应差异分析结果中，展示同一候选基因在单个细胞系 LRC vs BULK 比较中的方向、效应量和显著性。}
  {results/ngs/GSE114012/intersect/DLD1_HCT15_SW48/HCT15/deg_results.tex}


\FrameTable
  {02. DLD1/HCT15/SW48 方案中的 SW48 DEG 子集}
  {results/ngs/GSE114012/intersect/DLD1_HCT15_SW48/SW48/deg_results.csv}
  {该表将交集基因回填到对应差异分析结果中，展示同一候选基因在单个细胞系 LRC vs BULK 比较中的方向、效应量和显著性。}
  {results/ngs/GSE114012/intersect/DLD1_HCT15_SW48/SW48/deg_results.tex}


\FrameTable
  {02. HT55/SW948 交集统计}
  {results/ngs/GSE114012/intersect/HT55_SW948/summary.csv}
  {该表汇总当前交集方案中跨细胞系重复出现的显著基因数量及共同上调/下调方向，用于优先筛选稳定的休眠维持或复苏相关候选信号。}
  {results/ngs/GSE114012/intersect/HT55_SW948/summary.tex}


\FrameTable
  {02. HT55/SW948 交集基因列表}
  {results/ngs/GSE114012/intersect/HT55_SW948/gene_list.csv}
  {该表列出交集后的基因清单和基础注释，去除了模型统计量，便于将稳定候选基因继续输入 TF 富集、通路富集和药物重定位分析。}
  {results/ngs/GSE114012/intersect/HT55_SW948/gene_list.tex}


\FrameTable
  {02. HT55/SW948 方案中的 HT55 DEG 子集}
  {results/ngs/GSE114012/intersect/HT55_SW948/HT55/deg_results.csv}
  {该表将交集基因回填到对应差异分析结果中，展示同一候选基因在单个细胞系 LRC vs BULK 比较中的方向、效应量和显著性。}
  {results/ngs/GSE114012/intersect/HT55_SW948/HT55/deg_results.tex}


\FrameTable
  {02. HT55/SW948 方案中的 SW948 DEG 子集}
  {results/ngs/GSE114012/intersect/HT55_SW948/SW948/deg_results.csv}
  {该表将交集基因回填到对应差异分析结果中，展示同一候选基因在单个细胞系 LRC vs BULK 比较中的方向、效应量和显著性。}
  {results/ngs/GSE114012/intersect/HT55_SW948/SW948/deg_results.tex}


\FrameTable
  {02. SW948/RKO 交集统计}
  {results/ngs/GSE114012/intersect/SW948_RKO/summary.csv}
  {该表汇总当前交集方案中跨细胞系重复出现的显著基因数量及共同上调/下调方向，用于优先筛选稳定的休眠维持或复苏相关候选信号。}
  {results/ngs/GSE114012/intersect/SW948_RKO/summary.tex}


\FrameTable
  {02. SW948/RKO 交集基因列表}
  {results/ngs/GSE114012/intersect/SW948_RKO/gene_list.csv}
  {该表列出交集后的基因清单和基础注释，去除了模型统计量，便于将稳定候选基因继续输入 TF 富集、通路富集和药物重定位分析。}
  {results/ngs/GSE114012/intersect/SW948_RKO/gene_list.tex}


\FrameTable
  {02. SW948/RKO 方案中的 SW948 DEG 子集}
  {results/ngs/GSE114012/intersect/SW948_RKO/SW948/deg_results.csv}
  {该表将交集基因回填到对应差异分析结果中，展示同一候选基因在单个细胞系 LRC vs BULK 比较中的方向、效应量和显著性。}
  {results/ngs/GSE114012/intersect/SW948_RKO/SW948/deg_results.tex}


\FrameTable
  {02. SW948/RKO 方案中的 RKO DEG 子集}
  {results/ngs/GSE114012/intersect/SW948_RKO/RKO/deg_results.csv}
  {该表将交集基因回填到对应差异分析结果中，展示同一候选基因在单个细胞系 LRC vs BULK 比较中的方向、效应量和显著性。}
  {results/ngs/GSE114012/intersect/SW948_RKO/RKO/deg_results.tex}


\FrameFig
  {03. 传统火山图：ALL}
  {results/ngs/GSE114012/plots/volcano/ALL/volcano_plot.pdf}
  {单组火山图（ALL）}
  {
    \item 火山图展示 LRC 相对 BULK 的效应量与显著性分布，红/蓝分别提示休眠样状态中的上调和下调信号；
    \item 标注基因可作为后续 TF 调控、通路富集和药物逆转 signature 的候选入口；
    \item 单组图用于判断该模型中休眠-增殖状态切换的主要表达漂移方向。
  }

\FrameFig
  {03. 传统火山图：DLD1}
  {results/ngs/GSE114012/plots/volcano/DLD1/volcano_plot.pdf}
  {单组火山图（DLD1）}
  {
    \item 火山图展示 LRC 相对 BULK 的效应量与显著性分布，红/蓝分别提示休眠样状态中的上调和下调信号；
    \item 标注基因可作为后续 TF 调控、通路富集和药物逆转 signature 的候选入口；
    \item 单组图用于判断该模型中休眠-增殖状态切换的主要表达漂移方向。
  }

\FrameFig
  {03. 传统火山图：DLD1/HCT15}
  {results/ngs/GSE114012/plots/volcano/DLD1_HCT15/volcano_plot.pdf}
  {单组火山图（DLD1/HCT15）}
  {
    \item 火山图展示 LRC 相对 BULK 的效应量与显著性分布，红/蓝分别提示休眠样状态中的上调和下调信号；
    \item 标注基因可作为后续 TF 调控、通路富集和药物逆转 signature 的候选入口；
    \item 单组图用于判断该模型中休眠-增殖状态切换的主要表达漂移方向。
  }

\FrameFig
  {03. 传统火山图：DLD1/HCT15/SW48}
  {results/ngs/GSE114012/plots/volcano/DLD1_HCT15_SW48/volcano_plot.pdf}
  {单组火山图（DLD1/HCT15/SW48）}
  {
    \item 火山图展示 LRC 相对 BULK 的效应量与显著性分布，红/蓝分别提示休眠样状态中的上调和下调信号；
    \item 标注基因可作为后续 TF 调控、通路富集和药物逆转 signature 的候选入口；
    \item 单组图用于判断该模型中休眠-增殖状态切换的主要表达漂移方向。
  }

\FrameFig
  {03. 传统火山图：HCT15}
  {results/ngs/GSE114012/plots/volcano/HCT15/volcano_plot.pdf}
  {单组火山图（HCT15）}
  {
    \item 火山图展示 LRC 相对 BULK 的效应量与显著性分布，红/蓝分别提示休眠样状态中的上调和下调信号；
    \item 标注基因可作为后续 TF 调控、通路富集和药物逆转 signature 的候选入口；
    \item 单组图用于判断该模型中休眠-增殖状态切换的主要表达漂移方向。
  }

\FrameFig
  {03. 传统火山图：HT55}
  {results/ngs/GSE114012/plots/volcano/HT55/volcano_plot.pdf}
  {单组火山图（HT55）}
  {
    \item 火山图展示 LRC 相对 BULK 的效应量与显著性分布，红/蓝分别提示休眠样状态中的上调和下调信号；
    \item 标注基因可作为后续 TF 调控、通路富集和药物逆转 signature 的候选入口；
    \item 单组图用于判断该模型中休眠-增殖状态切换的主要表达漂移方向。
  }

\FrameFig
  {03. 传统火山图：RKO}
  {results/ngs/GSE114012/plots/volcano/RKO/volcano_plot.pdf}
  {单组火山图（RKO）}
  {
    \item 火山图展示 LRC 相对 BULK 的效应量与显著性分布，红/蓝分别提示休眠样状态中的上调和下调信号；
    \item 标注基因可作为后续 TF 调控、通路富集和药物逆转 signature 的候选入口；
    \item 单组图用于判断该模型中休眠-增殖状态切换的主要表达漂移方向。
  }

\FrameFig
  {03. 传统火山图：SW48}
  {results/ngs/GSE114012/plots/volcano/SW48/volcano_plot.pdf}
  {单组火山图（SW48）}
  {
    \item 火山图展示 LRC 相对 BULK 的效应量与显著性分布，红/蓝分别提示休眠样状态中的上调和下调信号；
    \item 标注基因可作为后续 TF 调控、通路富集和药物逆转 signature 的候选入口；
    \item 单组图用于判断该模型中休眠-增殖状态切换的主要表达漂移方向。
  }

\FrameFig
  {03. 传统火山图：SW948}
  {results/ngs/GSE114012/plots/volcano/SW948/volcano_plot.pdf}
  {单组火山图（SW948）}
  {
    \item 火山图展示 LRC 相对 BULK 的效应量与显著性分布，红/蓝分别提示休眠样状态中的上调和下调信号；
    \item 标注基因可作为后续 TF 调控、通路富集和药物逆转 signature 的候选入口；
    \item 单组图用于判断该模型中休眠-增殖状态切换的主要表达漂移方向。
  }

\FrameFig
  {04. 多组火山图：ALL}
  {results/ngs/GSE114012/plots/multiple_volcano/ALL/multiple_volcano_plot.pdf}
  {多组并列火山图（方案 ALL）}
  {
    \item 多组火山图并列展示多个细胞系或组合方案中的显著基因分布，便于观察跨模型方向一致性；
    \item 稳定重复出现的上调/下调信号更适合作为休眠维持或复苏触发轴的候选基因；
    \item 该图为从单模型差异走向共同调控网络提供直观筛选依据。
  }

\FrameFig
  {04. 多组火山图：DLD1/HCT15}
  {results/ngs/GSE114012/plots/multiple_volcano/DLD1_HCT15/multiple_volcano_plot.pdf}
  {多组并列火山图（方案 DLD1/HCT15）}
  {
    \item 多组火山图并列展示多个细胞系或组合方案中的显著基因分布，便于观察跨模型方向一致性；
    \item 稳定重复出现的上调/下调信号更适合作为休眠维持或复苏触发轴的候选基因；
    \item 该图为从单模型差异走向共同调控网络提供直观筛选依据。
  }

\FrameFig
  {04. 多组火山图：DLD1/HCT15/SW48}
  {results/ngs/GSE114012/plots/multiple_volcano/DLD1_HCT15_SW48/multiple_volcano_plot.pdf}
  {多组并列火山图（方案 DLD1/HCT15/SW48）}
  {
    \item 多组火山图并列展示多个细胞系或组合方案中的显著基因分布，便于观察跨模型方向一致性；
    \item 稳定重复出现的上调/下调信号更适合作为休眠维持或复苏触发轴的候选基因；
    \item 该图为从单模型差异走向共同调控网络提供直观筛选依据。
  }

\FrameFig
  {04. 多组火山图：HT55/SW948}
  {results/ngs/GSE114012/plots/multiple_volcano/HT55_SW948/multiple_volcano_plot.pdf}
  {多组并列火山图（方案 HT55/SW948）}
  {
    \item 多组火山图并列展示多个细胞系或组合方案中的显著基因分布，便于观察跨模型方向一致性；
    \item 稳定重复出现的上调/下调信号更适合作为休眠维持或复苏触发轴的候选基因；
    \item 该图为从单模型差异走向共同调控网络提供直观筛选依据。
  }

\FrameFig
  {04. 多组火山图：SW48/RKO}
  {results/ngs/GSE114012/plots/multiple_volcano/SW48_RKO/multiple_volcano_plot.pdf}
  {多组并列火山图（方案 SW48/RKO）}
  {
    \item 多组火山图并列展示多个细胞系或组合方案中的显著基因分布，便于观察跨模型方向一致性；
    \item 稳定重复出现的上调/下调信号更适合作为休眠维持或复苏触发轴的候选基因；
    \item 该图为从单模型差异走向共同调控网络提供直观筛选依据。
  }

\FrameFig
  {04. 多组火山图：COMBINATION 1}
  {results/ngs/GSE114012/plots/multiple_volcano/COMBINATION_1/multiple_volcano_plot.pdf}
  {多组并列火山图（方案 COMBINATION 1）}
  {
    \item 多组火山图并列展示多个细胞系或组合方案中的显著基因分布，便于观察跨模型方向一致性；
    \item 稳定重复出现的上调/下调信号更适合作为休眠维持或复苏触发轴的候选基因；
    \item 该图为从单模型差异走向共同调控网络提供直观筛选依据。
  }

\FrameFig
  {05. Top 50差异基因热图：ALL}
  {results/ngs/GSE114012/plots/gene_heatmap/ALL/gene_heatmap.pdf}
  {Top 50显著差异基因表达热图（ALL）}
  {
    \item 该图基于显著差异基因表按 P 值与 logFC 幅度筛选 Top 50 基因，并回到 TPM 矩阵中计算 \texttt{log2(TPM+1)} 行 Z-Score；
    \item LRC 统一置于左侧、BULK 置于右侧，便于直接判断休眠样状态与相对增殖状态之间的表达分离；
    \item 红蓝方向条带提示候选基因在 LRC 中的上调或下调方向，可作为后续休眠维持/复苏候选轴筛选依据。
  }

\FrameFig
  {05. Top 50差异基因热图：DLD1}
  {results/ngs/GSE114012/plots/gene_heatmap/DLD1/gene_heatmap.pdf}
  {Top 50显著差异基因表达热图（DLD1）}
  {
    \item 该图基于显著差异基因表按 P 值与 logFC 幅度筛选 Top 50 基因，并回到 TPM 矩阵中计算 \texttt{log2(TPM+1)} 行 Z-Score；
    \item LRC 统一置于左侧、BULK 置于右侧，便于直接判断休眠样状态与相对增殖状态之间的表达分离；
    \item 红蓝方向条带提示候选基因在 LRC 中的上调或下调方向，可作为后续休眠维持/复苏候选轴筛选依据。
  }

\FrameFig
  {05. Top 50差异基因热图：DLD1/HCT15}
  {results/ngs/GSE114012/plots/gene_heatmap/DLD1_HCT15/gene_heatmap.pdf}
  {Top 50显著差异基因表达热图（DLD1/HCT15）}
  {
    \item 该图基于显著差异基因表按 P 值与 logFC 幅度筛选 Top 50 基因，并回到 TPM 矩阵中计算 \texttt{log2(TPM+1)} 行 Z-Score；
    \item LRC 统一置于左侧、BULK 置于右侧，便于直接判断休眠样状态与相对增殖状态之间的表达分离；
    \item 跨细胞系组合更强调共同表达模式，有助于从单细胞系特异信号中提炼更稳健的休眠相关候选基因。
  }

\FrameFig
  {05. Top 50差异基因热图：DLD1/HCT15/SW48}
  {results/ngs/GSE114012/plots/gene_heatmap/DLD1_HCT15_SW48/gene_heatmap.pdf}
  {Top 50显著差异基因表达热图（DLD1/HCT15/SW48）}
  {
    \item 该图基于显著差异基因表按 P 值与 logFC 幅度筛选 Top 50 基因，并回到 TPM 矩阵中计算 \texttt{log2(TPM+1)} 行 Z-Score；
    \item LRC 统一置于左侧、BULK 置于右侧，便于直接判断休眠样状态与相对增殖状态之间的表达分离；
    \item 该组合可用于观察多模型共同休眠样表达特征，为后续 TF 调控、通路富集和药物重定位提供更严格入口。
  }

\FrameFig
  {05. Top 50差异基因热图：HCT15}
  {results/ngs/GSE114012/plots/gene_heatmap/HCT15/gene_heatmap.pdf}
  {Top 50显著差异基因表达热图（HCT15）}
  {
    \item 该图基于显著差异基因表按 P 值与 logFC 幅度筛选 Top 50 基因，并回到 TPM 矩阵中计算 \texttt{log2(TPM+1)} 行 Z-Score；
    \item LRC 统一置于左侧、BULK 置于右侧，便于直接判断休眠样状态与相对增殖状态之间的表达分离；
    \item 单细胞系热图可检查该模型自身的 LRC/BULK 表达分界，并辅助识别与组合分析一致的候选基因。
  }

\FrameFig
  {05. Top 50差异基因热图：HT55}
  {results/ngs/GSE114012/plots/gene_heatmap/HT55/gene_heatmap.pdf}
  {Top 50显著差异基因表达热图（HT55）}
  {
    \item 该图基于显著差异基因表按 P 值与 logFC 幅度筛选 Top 50 基因，并回到 TPM 矩阵中计算 \texttt{log2(TPM+1)} 行 Z-Score；
    \item LRC 统一置于左侧、BULK 置于右侧，便于直接判断休眠样状态与相对增殖状态之间的表达分离；
    \item 单细胞系热图可检查该模型自身的 LRC/BULK 表达分界，并辅助识别与组合分析一致的候选基因。
  }

\FrameFig
  {05. Top 50差异基因热图：RKO}
  {results/ngs/GSE114012/plots/gene_heatmap/RKO/gene_heatmap.pdf}
  {Top 50显著差异基因表达热图（RKO）}
  {
    \item 该图基于显著差异基因表按 P 值与 logFC 幅度筛选 Top 50 基因，并回到 TPM 矩阵中计算 \texttt{log2(TPM+1)} 行 Z-Score；
    \item LRC 统一置于左侧、BULK 置于右侧，便于直接判断休眠样状态与相对增殖状态之间的表达分离；
    \item 单细胞系热图可检查该模型自身的 LRC/BULK 表达分界，并辅助识别与组合分析一致的候选基因。
  }

\FrameFig
  {05. Top 50差异基因热图：SW48}
  {results/ngs/GSE114012/plots/gene_heatmap/SW48/gene_heatmap.pdf}
  {Top 50显著差异基因表达热图（SW48）}
  {
    \item 该图基于显著差异基因表按 P 值与 logFC 幅度筛选 Top 50 基因，并回到 TPM 矩阵中计算 \texttt{log2(TPM+1)} 行 Z-Score；
    \item LRC 统一置于左侧、BULK 置于右侧，便于直接判断休眠样状态与相对增殖状态之间的表达分离；
    \item 单细胞系热图可检查该模型自身的 LRC/BULK 表达分界，并辅助识别与组合分析一致的候选基因。
  }

\FrameFig
  {05. Top 50差异基因热图：SW948}
  {results/ngs/GSE114012/plots/gene_heatmap/SW948/gene_heatmap.pdf}
  {Top 50显著差异基因表达热图（SW948）}
  {
    \item 该图基于显著差异基因表按 P 值与 logFC 幅度筛选 Top 50 基因，并回到 TPM 矩阵中计算 \texttt{log2(TPM+1)} 行 Z-Score；
    \item LRC 统一置于左侧、BULK 置于右侧，便于直接判断休眠样状态与相对增殖状态之间的表达分离；
    \item 单细胞系热图可检查该模型自身的 LRC/BULK 表达分界，并辅助识别与组合分析一致的候选基因。
  }

\begin{frame}{06. GSEA分析说明}
  \begin{itemize}
    \item 06号脚本读取每套差异分析的 \texttt{all\_genes.csv}，以 limma 的 \texttt{t} 统计量构建排序基因列表，使用 \texttt{clusterProfiler::GSEA} 批量分析 msigdbr 当前数据库中的全部可用基因集类别。
    \item 当前全量运行覆盖 9 个差异分析设计和 26 类 MSigDB 基因集，共输出 234 套 GSEA 结果表、234 张 dotplot，以及 2057 张单通路 GSEA+核心基因表达热图。
    \item 结果中的 \texttt{core\_enrichment} 已转换为 Symbol，便于直接追踪富集到的基因；单通路热图使用 SE 对象 TPM，并以临床信息 \texttt{Title} 作为样本名。
    \item 该步骤用于从基因层面的 LRC/BULK 差异，上升到细胞周期、应激反应、代谢重编程、炎症信号和转录因子靶基因集等通路层面的休眠/复苏线索。
  \end{itemize}
\end{frame}

\FrameTable
  {06. GSEA全局Summary}
  {results/ngs/GSE114012/tables/GSEA_summary/summary.csv}
  {该表索引全部 GSEA 运行结果，统计每个分析设计、每类基因集的通路数量、NES方向和单通路图数量，是检查全量GSEA结果是否完整的总控表。}
  {results/ngs/GSE114012/tables/GSEA_summary/summary.tex}

\FrameFig
  {06. GSEA Dotplot：ALL / Hallmark}
  {results/ngs/GSE114012/plots/GSEA/ALL/hallmark/dotplot.pdf}
  {ALL设计的Hallmark GSEA气泡图}
  {
    \item Hallmark 基因集压缩了经典癌症与应激通路，适合快速判断 LRC 休眠样状态是否富集细胞周期、代谢、炎症或缺氧相关程序；
    \item 气泡大小表示核心富集比例，颜色反映显著性，正负NES方向帮助区分 LRC 相对 BULK 的激活或抑制趋势；
    \item ALL设计提供跨细胞系总体趋势，可作为后续聚焦通路的第一层筛选。
  }

\FrameFig
  {06. GSEA Dotplot：ALL / GO BP}
  {results/ngs/GSE114012/plots/GSEA/ALL/GO_BP/dotplot.pdf}
  {ALL设计的GO Biological Process GSEA气泡图}
  {
    \item GO BP 用更细粒度的生物过程描述 LRC/BULK 差异，有助于发现细胞周期重启、RNA加工、线粒体功能和应激反应等功能模块；
    \item 这些模块可与课题假说中的休眠维持和苏醒触发过程对应，作为后续 GSVA、TF 调控和药物重定位的候选通路输入；
    \item 通路名称在绘图中已精简前缀并自动换行，以减少长标签重叠。
  }

\FrameFig
  {06. GSEA Dotplot：DLD1/HCT15/SW48 / Hallmark}
  {results/ngs/GSE114012/plots/GSEA/DLD1_HCT15_SW48/hallmark/dotplot.pdf}
  {三细胞系组合设计的Hallmark GSEA气泡图}
  {
    \item DLD1/HCT15/SW48组合比单细胞系更严格，能突出跨模型重复出现的休眠样通路特征；
    \item 若细胞周期、MYC/E2F、炎症或缺氧相关通路在组合设计中仍保持显著，说明这些信号可能不是单一细胞系特异现象；
    \item 该图可与02号交集基因一起，用于筛选更稳健的候选调控轴。
  }

\FrameFig
  {06. GSEA Dotplot：DLD1/HCT15/SW48 / GO BP}
  {results/ngs/GSE114012/plots/GSEA/DLD1_HCT15_SW48/GO_BP/dotplot.pdf}
  {三细胞系组合设计的GO BP GSEA气泡图}
  {
    \item GO BP 结果用于解析组合模型中的共同生物过程变化，尤其关注增殖、染色体组织、RNA处理、代谢和应激相关过程；
    \item 与单基因火山图相比，GSEA可以保留未达到单基因阈值但整体方向一致的弱信号；
    \item 这些通路为后续“休眠细胞苏醒轨迹”和关键TF靶基因富集提供候选方向。
  }

\FrameFig
  {06. 单通路GSEA：ALL / Hypoxia}
  {results/ngs/GSE114012/plots/GSEA/ALL/hallmark/single_pathway/08_HALLMARK_HYPOXIA/gsea_plot.pdf}
  {ALL设计中Hypoxia通路的GSEA曲线与核心基因表达热图}
  {
    \item 单通路图上半部分展示基因集沿排序基因列表的富集曲线，下半部分展示核心富集基因在各样本中的 TPM 行 Z-score；
    \item 样本名使用临床 \texttt{Title}，便于直接分辨 LRC/BULK 和细胞系来源；
    \item 缺氧与代谢重编程常与休眠生态位和复苏触发有关，可作为后续机制验证的候选方向。
  }

\FrameFig
  {06. 单通路GSEA：ALL / TNFA-NFKB}
  {results/ngs/GSE114012/plots/GSEA/ALL/hallmark/single_pathway/06_HALLMARK_TNFA_SIGNALING_VIA_NFKB/gsea_plot.pdf}
  {ALL设计中TNFA/NFKB通路的GSEA曲线与核心基因表达热图}
  {
    \item TNF/NF-kappaB 轴与炎症反应、免疫微环境刺激和癌细胞状态切换密切相关；
    \item 若该通路在 LRC/BULK 排序中呈现稳定富集，可与课题假说中的 macrophage TNF/IL1beta 到 NF-kappaB 苏醒轴对应；
    \item 核心富集基因热图可帮助判断该通路信号是否在 LRC 样本中形成一致表达模式。
  }

\end{document}
